\ticket{9}{Лингвистические онтологии и тезаурусы}

\begin{examplan}
    \item Определить тезаурус и лингвистическую онтологию.
    \item Описать особенности их состава и построения.
    \item Перечислить виды лингвистических онтологических ресурсов.
    \item Раскрыть основные семантические отношения с примерами.
\end{examplan}

\subsection*{Короткая версия для письменного ответа}

\textbf{Тезаурус} --- словарь, в котором лексические единицы связаны семантическими отношениями: синонимией, антонимией, гипонимией, меронимией и ассоциациями. Он помогает навигации по смысловым связям слов.

\textbf{Лингвистическая онтология} --- онтологический ресурс, описывающий язык или языковые значения: слова, значения, грамматические категории, семантические роли и отношения между ними.

\subsection*{Тезаурусы}

Основные элементы тезауруса:
\begin{itemize}
    \item лексические единицы: слова и устойчивые словосочетания;
    \item синсеты --- группы слов с близким значением;
    \item краткие определения и примеры употребления;
    \item частеречная информация;
    \item семантические связи между словами или синсетами.
\end{itemize}

Примеры: WordNet, RuWordNet, RuThes, тезаурус Роже.

Виды тезаурусов:
\begin{itemize}
    \item идеографические, где слова организованы по смысловым областям;
    \item информационно-поисковые, используемые для индексирования и поиска;
    \item WordNet-подобные, подробно описывающие лексическую систему языка.
\end{itemize}

\subsection*{Лингвистические онтологии}

Лингвистическая онтология описывает не внешнюю предметную область, а языковую систему или связь языковых выражений с понятиями.

Она может включать:
\begin{itemize}
    \item лексические единицы, морфемы и фразы;
    \item грамматические категории: часть речи, падеж, число, время;
    \item синтаксические конструкции;
    \item семантические роли: агент, пациенс, инструмент;
    \item правила согласования, словообразования и интерпретации.
\end{itemize}

Особенность построения: часто используется путь \textbf{от слов к понятиям}. Анализируются реальные тексты и словоупотребления, затем выделяются значения и концепты. Поэтому перенос такой онтологии на другой язык требует адаптации: слова и границы значений в разных языках не совпадают полностью.

\subsection*{Виды лингвистических ресурсов}

\begin{itemize}
    \item \textbf{Однородные классификации}: ресурсы узкой области с одним типом отношений.
    \item \textbf{Тезаурусы}: понятия или значения связаны синонимией, иерархией и ассоциациями.
    \item \textbf{Рубрикаторы}: иерархические или фасетные классификации понятий и терминов.
    \item \textbf{Формальные лингвистические онтологии}: ресурсы в RDF/OWL с классами, свойствами и ограничениями.
\end{itemize}

\subsection*{Семантические отношения}

\begin{description}
    \item[Синонимия] близость значений: \textit{автомобиль, машина, авто}.
    \item[Антонимия] противоположность: \textit{большой -- маленький}.
    \item[Гипонимия/гиперонимия] отношение вид-род: \textit{пудель} --- гипоним слова \textit{собака}; \textit{собака} --- гипероним слова \textit{пудель}.
    \item[Меронимия/холонимия] отношение часть-целое: \textit{колесо} --- часть \textit{автомобиля}.
    \item[Тропонимия] для глаголов: более специфичный способ действия, например \textit{шептать} как способ \textit{говорить}.
    \item[Причинность и следование] одно событие вызывает или влечет другое: \textit{убить} влечет \textit{умереть}.
    \item[Ассоциация] смысловая или тематическая связь без строгого логического типа.
\end{description}

В формальных онтологиях отношения могут быть представлены через \texttt{rdfs:subClassOf}, \texttt{rdf:type}, объектные свойства, свойства данных, эквивалентность, несовместимость, транзитивность и симметричность.

\begin{important}
    \item Тезаурус фокусируется на словах и значениях, онтология --- на формальном описании концептов и отношений.
    \item Синсет объединяет лексические единицы с одним или близким значением.
    \item Лингвистические онтологии часто строятся от реального словоупотребления к понятиям.
    \item Главные отношения: синонимия, антонимия, гипонимия, меронимия, ассоциация.
    \item WordNet можно рассматривать как тезаурус и как легковесную лексическую онтологию.
\end{important}

